\documentclass[12pt,a4paper]{article}
\usepackage[utf8]{inputenc}
\usepackage[T1]{fontenc}
\usepackage[french]{babel}
\usepackage{geometry}
\usepackage{graphicx}
\usepackage{listings}
\usepackage{xcolor}
\usepackage{hyperref}
\usepackage{enumitem}
\usepackage{titlesec}
\usepackage{fancyhdr}
\usepackage{array}
\usepackage{booktabs}
\usepackage{longtable}

\geometry{margin=2.5cm}

\definecolor{codegreen}{rgb}{0,0.6,0}
\definecolor{codegray}{rgb}{0.5,0.5,0.5}
\definecolor{codepurple}{rgb}{0.58,0,0.82}
\definecolor{backcolour}{rgb}{0.95,0.95,0.95}

\lstdefinestyle{mystyle}{
    backgroundcolor=\color{backcolour},
    commentstyle=\color{codegreen},
    keywordstyle=\color{codepurple}\bfseries,
    stringstyle=\color{codegreen},
    basicstyle=\ttfamily\small,
    breaklines=true,
    frame=single,
    numbers=left,
    numberstyle=\tiny\color{codegray},
    tabsize=2
}
\lstset{style=mystyle}

\pagestyle{fancy}
\fancyhf{}
\rhead{Groupe 25 -- Master DI}
\lhead{TP Examen -- Gestion de Commandes}
\cfoot{\thepage}

\title{
    \vspace{-1cm}
    \textbf{Rapport TP Examen} \\
    \large Gestion d'une Ligne de Commande \\
    \large Application Android \\
    \vspace{0.5cm}
    \normalsize Master DI -- Programmation Web \\
    Universit\'e de Moundou
}
\author{
    \textbf{Groupe 25} \\
    Ngomde Cadet Kamdaou \\
    Bediang Leandrine Orline
}
\date{Octobre 2026}

\begin{document}

\maketitle
\thispagestyle{empty}

\newpage
\tableofcontents
\newpage

%=============================================================================
\section{Introduction}

Ce rapport pr\'esente la r\'ealisation du TP Examen de Programmation Web, portant sur la \textbf{gestion d'une ligne de commande}. Le projet a \'et\'e d\'evelopp\'e sous forme d'une application mobile Android native utilisant les technologies modernes de l'\'ecosyst\`eme Android.

L'application permet de g\'erer des produits, des commandes, des clients, et offre un tableau de bord statistique complet.

\subsection{Objectifs}
\begin{enumerate}
    \item Proposer un formulaire d'enregistrement des produits d'une commande donn\'ee
    \item Afficher la liste des produits, des commandes et des clients
    \item Modifier un produit
    \item Supprimer un produit dans une commande
    \item R\'ealiser un tableau de bord avec des statistiques sur les lignes de commande
\end{enumerate}

%=============================================================================
\section{Technologies et Architecture}

\subsection{Stack technique}

\begin{table}[h]
\centering
\begin{tabular}{ll}
\toprule
\textbf{Composant} & \textbf{Technologie} \\
\midrule
Langage & Kotlin 2.0 \\
Interface utilisateur & Jetpack Compose + Material 3 \\
Base de donn\'ees & Room (SQLite) \\
Injection de d\'ependances & Koin 3.5 \\
Navigation & Navigation Compose \\
Chargement d'images & Coil \\
Architecture & MVVM + Clean Architecture \\
Min SDK & API 24 (Android 7.0) \\
\bottomrule
\end{tabular}
\caption{Technologies utilis\'ees}
\end{table}

\subsection{Architecture Clean Architecture}

L'application suit le patron \textbf{Clean Architecture} avec trois couches bien distinctes par fonctionnalit\'e :

\begin{itemize}
    \item \textbf{Data} -- Contient les impl\'ementations des repositories, les acc\`es \`a la base de donn\'ees (DAOs, Entities)
    \item \textbf{Domain} -- Contient les mod\`eles m\'etier, les interfaces des repositories et les cas d'utilisation (use cases)
    \item \textbf{Presentation} -- Contient les ViewModels, les \'etats UI (State) et les \'ecrans Compose
\end{itemize}

Le patron \textbf{Repository} assure la s\'eparation entre la source de donn\'ees et la logique m\'etier. Chaque feature poss\`ede son propre ensemble de couches, favorisant la modularit\'e et la testabilit\'e.

\subsection{Pattern MVVM}

Chaque \'ecran suit le patron MVVM :
\begin{itemize}
    \item Le \textbf{ViewModel} expose un \texttt{StateFlow} repr\'esentant l'\'etat de l'UI
    \item L'\'ecran \textbf{Compose} observe cet \'etat via \texttt{collectAsStateWithLifecycle}
    \item Les actions utilisateur sont propag\'ees au ViewModel via des fonctions
\end{itemize}

\subsection{Internationalisation}

L'application supporte le \textbf{fran\c{c}ais} et l'\textbf{anglais} gr\^ace \`a :
\begin{itemize}
    \item Des fichiers \texttt{strings.xml} dans \texttt{values/} (anglais) et \texttt{values-fr/} (fran\c{c}ais)
    \item Une classe \texttt{UiText} scell\'ee permettant de manipuler les textes de mani\`ere type-safe
\end{itemize}

%=============================================================================
\section{Mod\`ele de donn\'ees}

\subsection{Sch\'ema de la base de donn\'ees}

La base de donn\'ees comporte les tables suivantes :

\begin{itemize}
    \item \texttt{Client(idCli, nomCli, prenomCli, email, photo, idRegion)}
    \item \texttt{Commande(idCmd, dateCmd, \#idCli)}
    \item \texttt{ligneCommande(\#idCmd, \#idProd, qteCmd)} -- cl\'e primaire composite
    \item \texttt{Produit(idProd, designProd, qteStock, prixUProd, photoProd)}
    \item \texttt{Region(idRegion, nomRegion, \#idDep)}
    \item \texttt{D\'epartement(idDep, nomDep)}
    \item \texttt{Users(id, login, password)} -- pour l'authentification
\end{itemize}

\subsection{Relations}

\begin{itemize}
    \item Un \textbf{D\'epartement} contient plusieurs \textbf{R\'egions}
    \item Une \textbf{R\'egion} contient plusieurs \textbf{Clients}
    \item Un \textbf{Client} peut passer plusieurs \textbf{Commandes}
    \item Une \textbf{Commande} contient plusieurs \textbf{Produits} via \texttt{ligneCommande}
    \item Un \textbf{Produit} peut appara\^itre dans plusieurs \textbf{Commandes}
\end{itemize}

Les contraintes d'int\'egrit\'e r\'ef\'erentielle sont appliqu\'ees via des \texttt{ForeignKey} Room avec \texttt{CASCADE} sur la suppression.

\subsection{Pr\'e-remplissage}

La base est pr\'e-remplie au premier lancement avec :
\begin{itemize}
    \item 1 utilisateur pour l'authentification
    \item 3 d\'epartements, 5 r\'egions
    \item 8 clients
    \item 12 produits
    \item 10 commandes avec 22 lignes de commande
\end{itemize}

%=============================================================================
\section{Fonctionnalit\'es r\'ealis\'ees}

\subsection{Authentification (Feature 0)}

Un \'ecran de connexion permet \`a l'utilisateur de s'authentifier avec les identifiants fournis :
\begin{itemize}
    \item \textbf{Login} : \texttt{MasTMound2}
    \item \textbf{Mot de passe} : \texttt{MoundMasT2@2026}
\end{itemize}

Fonctionnalit\'es :
\begin{itemize}
    \item Validation des champs (champs vides, identifiants incorrects)
    \item Masquage/affichage du mot de passe
    \item Messages d'erreur bilingues via \texttt{UiText}
    \item Redirection automatique apr\`es connexion r\'eussie
\end{itemize}

\subsection{Gestion des Produits (Features 2 et 3)}

\subsubsection{Liste des produits}
\begin{itemize}
    \item Affichage en liste avec photo, d\'esignation, prix et stock
    \item Barre de recherche en temps r\'eel (filtre par d\'esignation)
    \item \'Etat vide avec ic\^one quand aucun produit
\end{itemize}

\subsubsection{Ajout / Modification d'un produit}
\begin{itemize}
    \item Formulaire avec : d\'esignation, quantit\'e en stock, prix unitaire
    \item S\'election de photo depuis la galerie (stock\'ee en fichier local)
    \item Validation des champs avec messages d'erreur
    \item M\^eme \'ecran pour l'ajout et la modification (d\'etection via param\`etre de navigation)
\end{itemize}

\subsubsection{Suppression d'un produit}
\begin{itemize}
    \item Bo\^ite de dialogue de confirmation avant suppression
    \item Suppression en cascade des lignes de commande associ\'ees
\end{itemize}

\subsection{Gestion des Commandes (Features 1 et 4)}

\subsubsection{Liste des commandes}
\begin{itemize}
    \item Affichage avec num\'ero, date, nom du client, nombre d'articles et total
    \item Calcul automatique du montant total via jointure SQL
    \item Suppression avec confirmation
\end{itemize}

\subsubsection{Cr\'eation d'une commande}
\begin{itemize}
    \item S\'election du client via menu d\'eroulant
    \item Date automatique (timestamp actuel)
    \item Redirection vers le d\'etail apr\`es cr\'eation
\end{itemize}

\subsubsection{D\'etail d'une commande -- Enregistrement des produits}
\begin{itemize}
    \item Affichage des lignes de commande (produit, quantit\'e, prix, sous-total)
    \item Ajout d'un produit via bo\^ite de dialogue (s\'election produit + quantit\'e)
    \item Suppression d'un produit de la commande
    \item Affichage du montant total de la commande
\end{itemize}

\subsection{Liste des Clients (Feature 2)}

\begin{itemize}
    \item Affichage en cartes avec photo, nom complet, email et r\'egion
    \item Jointure SQL avec la table \texttt{region} pour afficher le nom de la r\'egion
    \item Donn\'ees issues du pr\'e-remplissage
\end{itemize}

\subsection{Tableau de Bord (Feature 5)}

Le tableau de bord offre une vue analytique compl\`ete avec les statistiques suivantes :

\subsubsection{Indicateurs cl\'es (KPI)}
\begin{itemize}
    \item \textbf{Chiffre d'affaires total} -- somme de (quantit\'e $\times$ prix unitaire) sur toutes les lignes
    \item \textbf{Nombre total de commandes}
    \item \textbf{Nombre total de produits}
    \item \textbf{Nombre total de clients}
\end{itemize}

\subsubsection{Top 5 des produits les plus vendus}
Requ\^ete agr\'eg\'ee par produit, tri\'ee par quantit\'e totale vendue, avec le chiffre d'affaires g\'en\'er\'e par chaque produit.

\subsubsection{Chiffre d'affaires par r\'egion}
Visualisation sous forme de barres horizontales proportionnelles, calculant le revenu total par r\'egion \`a travers la cha\^ine : R\'egion $\rightarrow$ Client $\rightarrow$ Commande $\rightarrow$ Ligne de commande $\rightarrow$ Produit.

\subsubsection{Commandes par client}
Tableau montrant pour chaque client le nombre de commandes et le montant total d\'epens\'e.

\subsubsection{Performance mensuelle}
Regroupement des commandes par mois avec le nombre de commandes et le revenu g\'en\'er\'e, sous forme de barres horizontales.

\subsubsection{Justification des choix}
Le tableau de bord a \'et\'e con\c{c}u pour offrir une vision \`a plusieurs niveaux :
\begin{itemize}
    \item \textbf{Vue macro} : les KPI donnent un aper\c{c}u instantan\'e de l'activit\'e
    \item \textbf{Vue produit} : identifier les best-sellers pour optimiser le stock
    \item \textbf{Vue g\'eographique} : comprendre la r\'epartition du CA par r\'egion
    \item \textbf{Vue client} : identifier les clients les plus actifs
    \item \textbf{Vue temporelle} : d\'etecter les tendances et la saisonnalit\'e
\end{itemize}

Les requ\^etes SQL sont optimis\'ees avec des jointures et agr\'egations directement au niveau de la base, \'evitant le traitement c\^ot\'e applicatif.

%=============================================================================
\section{Structure du Projet}

\begin{lstlisting}[language={}, title={Arborescence des packages}]
com.groupe25.tpexam3/
|-- core/
|   |-- data/local/
|   |   |-- dao/         (UserDao, ProductDao, ClientDao,
|   |   |                 CommandeDao, LigneCommandeDao,
|   |   |                 DashboardDao)
|   |   |-- entity/      (7 entites Room)
|   |   |-- AppDatabase.kt
|   |-- presentation/
|       |-- UiText.kt
|-- di/
|   |-- AppModule.kt     (configuration Koin)
|-- feature/
|   |-- auth/
|   |   |-- data/repository/AuthRepositoryImpl.kt
|   |   |-- domain/model/User.kt
|   |   |-- domain/repository/AuthRepository.kt
|   |   |-- domain/usecase/LoginUseCase.kt
|   |   |-- presentation/LoginScreen.kt
|   |   |-- presentation/LoginViewModel.kt
|   |-- product/
|   |   |-- data/repository/ProductRepositoryImpl.kt
|   |   |-- domain/model/Product.kt
|   |   |-- domain/repository/ProductRepository.kt
|   |   |-- presentation/list/ProductListScreen.kt
|   |   |-- presentation/addedit/AddEditProductScreen.kt
|   |-- order/
|   |   |-- data/repository/OrderRepositoryImpl.kt
|   |   |-- domain/model/Order.kt, OrderLine.kt
|   |   |-- domain/repository/OrderRepository.kt
|   |   |-- presentation/list/OrderListScreen.kt
|   |   |-- presentation/detail/OrderDetailScreen.kt
|   |   |-- presentation/create/CreateOrderScreen.kt
|   |-- client/
|   |   |-- presentation/ClientListScreen.kt
|   |-- dashboard/
|       |-- presentation/DashboardScreen.kt
|-- navigation/
|   |-- Screen.kt, BottomNavItem
|-- ui/theme/Theme.kt
|-- MainActivity.kt
|-- TpExamApp.kt
\end{lstlisting}

%=============================================================================
\section{Injection de d\'ependances}

Koin est configur\'e dans un module unique \texttt{AppModule.kt} qui fournit :
\begin{itemize}
    \item L'instance singleton de la base de donn\'ees \texttt{AppDatabase}
    \item Tous les DAOs en tant que singletons
    \item Les repositories en tant que singletons
    \item Les use cases en tant que factory
    \item Les ViewModels avec \texttt{viewModel \{ \}}
\end{itemize}

%=============================================================================
\section{Navigation}

L'application utilise une \textbf{Bottom Navigation Bar} avec 4 onglets :
\begin{enumerate}
    \item \textbf{Produits} -- Liste et CRUD des produits
    \item \textbf{Commandes} -- Liste, cr\'eation et d\'etail des commandes
    \item \textbf{Clients} -- Liste des clients
    \item \textbf{Tableau de bord} -- Statistiques et d\'econnexion
\end{enumerate}

Les \'ecrans secondaires (ajout/modification produit, d\'etail commande, cr\'eation commande) s'ouvrent sans la barre de navigation pour une exp\'erience utilisateur fluide.

%=============================================================================
\section{Conclusion}

Ce projet a permis de d\'evelopper une application mobile Android compl\`ete r\'epondant \`a l'ensemble des exigences du TP Examen. L'utilisation de Clean Architecture avec le pattern Repository assure une s\'eparation claire des responsabilit\'es, facilitant la maintenance et l'\'evolution du code.

Les points forts de la r\'ealisation :
\begin{itemize}
    \item Architecture modulaire et extensible
    \item Interface utilisateur moderne avec Material 3
    \item Support bilingue (fran\c{c}ais / anglais)
    \item Tableau de bord analytique riche avec des requ\^etes SQL optimis\'ees
    \item Gestion compl\`ete du cycle de vie des donn\'ees (CRUD)
    \item Pr\'e-remplissage de la base pour une d\'emonstration imm\'ediate
\end{itemize}

\end{document}
